\paragraph{AQ-WEIL-NOTATION: the recovered $Q'$ and the norm-one subgroup.}
The official cote 15 PDF has now been inspected beyond the
supplied excerpts.  Physical page 94, printed footer 93,
opens with the explicit definition
\emph{``$Q'$: cl\^oture tot[alement] r\'eelle de $\mathbb Q$''};
the bracket expands the handwritten abbreviation.
Thus $Q'=\mathbb Q^{\rm tr}$, the field of totally real
algebraic numbers, and the $i$ in $Q'(i)$ means
$\sqrt{-1}$, distinct from the integral weight index.
Physical page 98/footer 97 discusses a totally imaginary
quadratic extension of a totally real field and writes
$\lambda\overline\lambda=q^w$, followed by the
normalization $\lambda=q^{w/2}u$ with
$u\overline u=1$.
Physical page 100/footer 99 contains the boxed correction
requiring every conjugate to have absolute value $q^{w/2}$,
as already preserved in the supplied transcription.
These recovered norm and purity conditions identify the
weight-zero interpretation used in AQ-WEIL-LIMIT.

More precisely, let
\[
 \mathbb Q^{\rm CM}=\bigcup_{L\ {\mathrm{CM}}}L
 \quad\hbox{inside the fixed }\overline{\mathbb Q}\subset\mathbb C.
\]
Then
\[
 \mathbb Q^{\rm CM}=\mathbb Q^{\rm tr}(\sqrt{-1}).
 \tag{AQ.WN.1}
\]
For a prime power $q=p^n$, the group of algebraic
$p$-units with all conjugate absolute values $q^{w/2}$,
for some integer $w$, is exactly
\[
 \left\{\lambda\in\mathbb Q^{\rm tr}(\sqrt{-1})^*:
   \lambda\text{ is a $p$-unit},\quad
   \lambda\overline\lambda=q^w\text{ for some }w\in\mathbb Z
 \right\}.                                                   \tag{AQ.WN.2}
\]
Its weight-zero subgroup is $W_0(p)$, independent of $n$.
Thus the historical group $M(q)$ and its weight-zero part
agree with $W(q)$ and $W_0(p)$ when the explicit
norm/purity conditions on pages 98 and 100 are retained.

\begin{proof}
For a CM field $L$, let $L^+$ be its totally real
subfield.  Choose $b\in L$ with
$\delta=b-\overline b\ne0$.  Then
$\delta^2\in L^+$, and $a=-\delta^2$ is totally positive:
every embedding sends $\delta$ to a nonzero purely
imaginary number.  The algebraic number $\sqrt a$
is totally real, because all its conjugates are real
square roots of positive conjugates of $a$.
Hence
$L=L^+(\delta)\subset\mathbb Q^{\rm tr}(\sqrt{-1})$.
Conversely any element of the latter field belongs to
$F(\sqrt{-1})$ for a finite totally real field $F$,
and $F(\sqrt{-1})$ is a CM field.  This proves (AQ.WN.1),
including the assertion that the displayed union is a field.

Let $\lambda$ have the absolute values in (AQ.WN.2),
and let $L$ be the Galois closure of $\mathbb Q(\lambda)$.
Write $c$ for complex conjugation on $L$.
For every $\sigma\in\operatorname{Gal}(L/\mathbb Q)$,
\[
 c(\sigma\lambda)=q^w/(\sigma\lambda).
\]
For every $\tau$ in the same Galois group,
$c\tau$ and $\tau c$ consequently agree on all the
conjugates $\sigma\lambda$, which generate $L$.
Thus $c$ is central.
If $c=1$, the Galois field $L$ is totally real.
Otherwise $L^+=L^{\langle c\rangle}$ is totally real,
$[L:L^+]=2$, and $L$ has no real embedding, so $L$ is CM.
In either case (AQ.WN.1) places $\lambda$ in
$\mathbb Q^{\rm tr}(\sqrt{-1})$; the norm identity follows
from its specified absolute value at the fixed embedding.

Conversely this field is Galois over $\mathbb Q$, and
complex conjugation fixes $\mathbb Q^{\rm tr}$ and
negates $\sqrt{-1}$.  It therefore commutes with every
embedding.  Applying any embedding to
$\lambda\overline\lambda=q^w$ gives
$|\sigma\lambda|^2=q^w$, as required.
The $p$-unit condition is the same on both sides,
so (AQ.WN.2) follows.  Setting $w=0$ gives the claimed
weight-zero subgroup, independently of $q=p^n$.
\end{proof}

Finite-prime valuations alone must not be substituted
for the norm/purity condition in this identification.
For example $\varepsilon=2+\sqrt3$ and its inverse
$2-\sqrt3$ are algebraic integers, so $\varepsilon$
is a unit at every finite place.  In particular every
finite-prime sum $v_{\mathfrak p}(\varepsilon)+
v_{\overline{\mathfrak p}}(\varepsilon)$ is zero,
for every chosen $p$.  But
$\varepsilon\overline\varepsilon=\varepsilon^2\ne1$
and its two real conjugates have unequal absolute values.
No positive power is a weight-zero Weil number.
This checks the limitation of a \emph{valuation-only}
paraphrase of the margin on 15/102; it does not assert
that the manuscript's terminology \emph{unit\'e absolue}
itself discards archimedean conditions.
The corrected group interpretation uses the explicit
norm condition actually present on the preceding sheet.
